\documentclass[a4paper, amsfonts, amssymb, amsmath, reprint, showkeys, nofootinbib, twoside]{revtex4-2}
\usepackage[english]{babel}
\usepackage[utf8]{inputenc}
\usepackage[colorinlistoftodos, color=green!40, prependcaption]{todonotes}
\input{preamble}
\usepackage[pdftex, pdftitle={Article}, pdfauthor={Author}]{hyperref} % Para enlaces en el PDF
%\setlength{\marginparwidth}{2.5cm}
%\bibliographystyle{apsrev4-1} 
\begin{document}
\title{Quantum walks for 1D anyons in a cavity}

\author{Ronaldo Navarro Ambriz}
    \email[Correspondence email address: ]{ron@estudiantes.unam.fisi}% Correo asociado 
    \affiliation{Instituto de Física, Universidad Nacional Autónoma de México, Ciudad de México, México} 
\author{Santiago Francisco Caballero Benítez}
    \email[Correspondence email address: ]{scaballero@fisica.unam.mx}% Correo asociado 
    \affiliation{Instituto de Física, Universidad Nacional Autónoma de México, Ciudad de México, México}

\date{\today} % Se puede cambiar o quitar, se usara para parte de llevar un registro 

\begin{abstract}
\lipsum[1]
\end{abstract}

\keywords{anyons, quantum matter, ultracold atoms}


\maketitle


\section{Introduction} 
Los anyones son ... 
\section{Main Idea} 
Del Hamiltoniano de bosones condicionados: 
\begin{equation}\label{HamideAny}
\mathcal{H}^b=-J\sum_{i=1}^M \left( \hat{b}_{i}^\dagger   \hat{b}_{i+1} e^{\textit{i} \theta \hat{n}_i} + h.c. \right) + \frac{U}{2} \sum_{i=1}^M \hat{n}_i \left( \hat{n}_i - 1 \right),
\end{equation}
Añadiendo el término de cavidad: 
\begin{equation}
\mathcal{H}_{cav}= g_{cav} \sum_{i,j=1}^M \left( -1 \right)^{i-j} \hat{n}_i \hat{n}_j = g_{cav} \left( \sum_{i=1}^M (-1)^i \hat{n}_i \right)^2, 
\end{equation} 
Obteniendo el Hamiltoniano extendido: 
\begin{equation} 
\mathcal{H}^b_{ext} = \mathcal{H}^b +\mathcal{H}_{cav}
\end{equation} 
The original Hamiltonian \eqref{HamideAny} breaks the inversion symmetry $\mathcal{I}$, as the particles in the chain adquire a phase due a site occupation. 
Result, only for a change of sign in both term we will observe a total reflection in the quantum walk
\begin{equation} 
\langle \hat{n}_j (t) \rangle_{+U, +g}  = \langle \hat{n}_j'´ (t) \rangle_{-U, -g} 
\end{equation} 
\begin{equation} 
\langle \hat{n}_j (t) \rangle_{+U, -g}  = \langle \hat{n}_j' (t) \rangle_{-U, +g} 
\end{equation}

\section{Conclusion}

\bibliographystyle{unsrt}
\bibliography{referencias}

\appendix*
\section{Appendix} \label{sec:appendix}
Del operador de paridad: 
\begin{equation} 
\mathcal{P} = e^{i\pi \sum_{r} \hat{n}_{2r+1}}    
\end{equation}

\end{document}
